\documentclass[10pt]{article}
\usepackage[margin=0.7in]{geometry}
\usepackage[T1]{fontenc}
\usepackage{lmodern}
\usepackage{amsmath,amssymb}
\usepackage{enumitem}
\setlist[enumerate]{leftmargin=*, itemsep=7mm}
\setlist[enumerate,2]{label=(\Alph*), leftmargin=2em, itemsep=0.05em, topsep=0.25em}
\pagestyle{plain}

\begin{document}
\begin{center}
  {\large\bfseries Video 6 Quiz: Real Numbers and Continuity}\\[0.25em]
  \textit{Choose the best answer for each question.}
\end{center}
\noindent Name:\ \rule{3.2in}{0.4pt}\hfill Date:\ \rule{1.25in}{0.4pt}
\vspace{0.7em}
\begin{enumerate}
  \item What statement in this video will be proven on the board?
  \begin{enumerate}
    \item A bounded non-decreasing sequence in $\mathbb{R}$ converges to its maximum.
    \item A bounded non-decreasing sequence in $\mathbb{R}$ contains a subsequence that converges.
    \item A bounded non-decreasing sequence in $\mathbb{R}$ has a maximum.
    \item A bounded non-decreasing sequence in $\mathbb{R}$ converges to its supremum.
  \end{enumerate}

  \item How does the video say uniform continuity strengthens continuity at a point?
  \begin{enumerate}
    \item It requires the same function value at every point.
    \item It allows $\delta$ to be chosen independently of $x_0$.
    \item It requires the domain to be a closed interval.
    \item It replaces $\epsilon$ with a fixed Lipschitz constant.
  \end{enumerate}

  \item What behavior at $x=0$ does the video use to explain why $f(x)=\sqrt{x}$ is not Lipschitz continuous on $[0,1]$?
  \begin{enumerate}
    \item Its derivative becomes infinite.
    \item Its value becomes negative.
    \item Its graph has a horizontal tangent.
    \item Its range is unbounded.
  \end{enumerate}

  \item At the end of the square-root example, what does the video ask students to pause and decide?
  \begin{enumerate}
    \item Whether $\sqrt{x}$ is differentiable at $0$.
    \item Whether $\sqrt{x}$ is uniformly continuous.
    \item Whether $\sqrt{x}$ has a maximum on $[0,1]$.
    \item Whether $\sqrt{x}$ maps $[0,1]$ onto $\mathbb{R}$.
  \end{enumerate}
\end{enumerate}
\end{document}
